
\documentclass[10pt,letterpaper]{article}

\usepackage[margin=0.82in]{geometry}
\usepackage{times}
\usepackage{microtype}
\usepackage{enumitem}
\usepackage{amsmath}
\usepackage{booktabs}
\usepackage{array}

\setlength{\parindent}{0pt}
\setlength{\parskip}{4pt}
\setlist{nosep,leftmargin=1.4em}
\setcounter{secnumdepth}{2}

\title{\vspace{-1.5em}\textbf{SPECTRE: An Event-Driven Vision-Language\\
Tactical Coach for Apex Legends}}
\author{AI Systems Architecture --- Midterm Progress Report\\
Utah State University}
\date{}

\begin{document}
\maketitle
\vspace{-2em}

\section{Introduction}

Competitive gameplay analysis is fundamentally a problem of contextual
reasoning. In Apex Legends, the significance of an action depends on
the player's position, available resources, teammate status, enemy
threats, and the events that immediately precede or follow it.
Performance statistics can summarize an engagement, but they seldom
explain whether a different decision would have produced a better
outcome.

SPECTRE addresses this limitation by combining structured visual
perception with vision-language reasoning. Rather than requiring a
general-purpose VLM to interpret every element of a gameplay frame,
the system extracts selected HUD information using conventional
computer vision and provides that information as temporal context
during multimodal analysis.

The project began with an architectural question: \textit{Can
structured gameplay observations improve the reliability and
usefulness of VLM-generated tactical feedback?}

At the midterm milestone, SPECTRE has progressed from a proposed
architecture to a functional offline prototype. The implemented
pipeline processes gameplay recordings, identifies screen states,
recognizes squad legends and equipped weapons, estimates shield
strength, constructs timestamped observations, and performs
sequential VLM analysis. The model produces structured descriptions
of gameplay events, potential mistakes, recommendations, and
uncertainties.

Initial experiments confirm that the components can operate together.
They also expose important limitations. The VLM sometimes
misinterprets ammunition mechanics, treats unavailable measurements
as evidence of player incapacitation, and fails to reconstruct
complete engagement outcomes. These findings provide a concrete
basis for revising the architecture toward stronger temporal
reasoning and game-specific knowledge.

\section{System Architecture}

SPECTRE uses a two-stage processing architecture that separates
visual measurement from multimodal interpretation.

\subsection{Offline Perception Pipeline}

The first stage, implemented in \texttt{extract.py}, processes
recorded gameplay using OpenCV. Source recordings have a resolution
of $3840 \times 1600$ pixels. Each sampled frame is evaluated using
screen-specific detectors configured through a common JSON file.

The perception stage produces a chronological match record
containing frame indices, timestamps, and extracted HUD observations.
Its output is independent of the VLM, allowing individual detectors
to be tested without repeating multimodal inference.

\subsection{Vision-Language Analysis Pipeline}

The second stage, implemented in \texttt{analyze.py}, reads the
previously generated match record and retrieves the corresponding
images directly from the original recording.

Frames and their associated HUD observations are combined into
chronological batches and submitted to a locally hosted Qwen3.5
vision-language model through the \texttt{llama.cpp} multimodal
interface.

The model is instructed to identify significant events, analyze
positioning and coordination, and propose improvements grounded
in visible evidence. Its output is serialized into a structured
JSON report.

The separation between these stages offers two practical benefits.
First, the expensive VLM need not repeat deterministic HUD
recognition. Second, the same extracted match data can support
different prompts, models, and analysis policies without
rerunning the perception pipeline.

\section{Implementation}

\subsection{HUD Configuration and Preprocessing}

A major implementation challenge is the complexity of Apex
Legends' interface. Relevant information appears in small
screen regions and may change appearance with shield type,
weapon selection, player status, or the currently open menu.

SPECTRE addresses this through a preprocessing notebook that
records regions of interest (ROIs) and visual references in
\texttt{hud\_regions.json}. The configuration separates
screen-specific extraction parameters from the runtime logic.

This calibration process is necessary because a detector that
performs well on normal gameplay may produce meaningless
results when the same screen coordinates contain a death-box
menu, inventory panel, or map overlay.

\subsection{Screen-State Classification}

The system recognizes five explicit non-gameplay screen states:
inventory, map, spectating, death-box interaction, and match
summary. Each is associated with a distinctive interface
template and a fixed-position ROI.

Normal gameplay is identified when these explicit interfaces
are not detected and the player's legend portrait provides
sufficient evidence of an active gameplay HUD.

A conservative classification policy is applied. If no
candidate satisfies the required similarity threshold, or
multiple candidates are insufficiently distinguishable,
the screen is classified as \texttt{unknown}.

Gameplay-specific extraction is performed only for frames
classified as normal gameplay.

This mechanism resolves a failure observed during development:
before screen classification was introduced, the shield detector
occasionally interpreted death-box interface elements as
equipped shields. The updated pipeline avoids applying that
detector to recognized death-box screens.

\subsection{Legend Recognition}

Legend recognition uses position-specific template matching.
Separate reference collections are maintained for the local
player and two teammates because portraits may have different
sizes or appearances across the three HUD positions.

The current implementation contains 28 reference portraits
for each position. Templates are loaded once, converted to
grayscale, and compared against the corresponding frame ROI
using normalized cross-correlation.

For candidate template $T$ and observed region $I$, the
detector selects

\[
    \hat{\ell} =
    \arg\max_{\ell}
    \operatorname{NCC}(I,T_{\ell}),
\]

subject to a minimum similarity threshold.

In an inspected sequence, the detector repeatedly identified
Wattson, Catalyst, and Conduit in their correct squad positions.
Accepted similarity scores commonly exceeded 0.93, with
several observations above 0.98.

These scores demonstrate strong template correspondence in
the tested footage; they should not be interpreted as
dataset-wide recognition accuracy.

\subsection{Shield Classification and Estimation}

Shield extraction combines calibrated color references with
horizontal fill measurement. The current configuration includes
six visual categories: white, blue, and purple shields, each
with a corresponding Amped variant.

Because the three squad positions use different screen
coordinates, separate reference regions are recorded for
each player.

The runtime detector evaluates candidate shield colors and
tracks matching pixels from the left edge of the shield bar
toward the right. Small interruptions are tolerated to
account for segmented visual elements.

The estimated shield strength is

\[
    \hat{s} =
    \frac{w_{\mathrm{filled}}}{w_{\mathrm{reference}}}
    s_{\max},
\]

where $s_{\max}$ denotes the nominal capacity associated
with the recognized shield type.

The left-anchored approach reduces the influence of isolated
matching pixels outside the filled region. Nevertheless,
visual effects, bar boundaries, and changes in shield
appearance can still affect the estimate. Amped-shield
classification remains limited when a partially depleted
Amped shield is visually indistinguishable from a normal one.

Shield values are therefore treated as approximate
observations rather than verified ground truth.

\subsection{Weapon Recognition}

Equipped weapon names are extracted from the corresponding
HUD region using Tesseract OCR and a restricted vocabulary
of Apex Legends weapons.

The initial implementation also attempted to recognize
magazine and reserve ammunition and associate the observed
values with the active weapon.

These experiments revealed two systematic errors. First,
dimmed ammunition placeholders were occasionally recognized
as numerical values. Second, the active weapon could be
identified incorrectly, causing ammunition readings to
be associated with the wrong weapon.

For the midterm prototype, the reliable component---weapon
identity---has been retained. Ammunition extraction is
excluded from the primary structured observations until
its validity can be established.

This decision reflects an important design principle:
an unavailable measurement is preferable to an incorrect
measurement presented as reliable context.

\subsection{Temporal Match Representation}

Each sampled frame contributes an observation containing
its timestamp, source frame index, recognized screen state,
and available HUD information.

The observations are stored in a match-level JSON timeline.
This structure preserves the sequence of changing gameplay
conditions and supports direct alignment with the original
recording.

An important distinction is maintained between measured
information and missing information. A value of
\texttt{null} indicates that the field was not observed;
it does not imply that the associated resource is empty
or that a player has been knocked down.

The current system records observations chronologically,
but a more comprehensive state tracker---including
last-known values, validity timestamps, and explicit
state transitions---remains under development.

\section{Experimental Evaluation}

\subsection{Experimental Setup}

The prototype was evaluated using recorded Apex Legends
gameplay and a locally hosted Qwen3.5 vision-language
model.

The initial configuration samples gameplay at one frame
per second and submits batches of ten chronologically
ordered frames. Each image is accompanied by its
corresponding structured HUD observation.

The VLM returns a JSON object containing a segment summary,
timestamped events, potential mistakes, suggested
improvements, and unresolved uncertainties.

Streaming output is enabled to observe generation behavior,
while execution logs retain screen classifications,
detector scores, and model token usage.

The present results constitute functional and qualitative
evaluation on selected gameplay segments rather than a
comprehensive accuracy benchmark.

\subsection{Non-Combat Sequence}

The first experiment analyzed a ten-second segment
covering movement, looting, and inventory interaction.

The model successfully recognized the player's movement
from an elevated position toward cover, interaction with
loot, opening of the inventory screen, and subsequent
repositioning. It also acknowledged that no visible
combat occurred during the analyzed interval.

The experiment established that structured observations
and their corresponding images could be combined into
a coherent chronological analysis.

One measured request contained approximately 17,600
prompt tokens, with approximately 20,800 total tokens
consumed. This illustrates the substantial computational
cost of multimodal interpretation even for a short
sequence.

However, the model also produced inaccurate ammunition
recommendations, demonstrating that successful visual
recognition does not necessarily imply correct knowledge
of game mechanics.

\subsection{Combat Sequence}

A second experiment examined the interval from
540 to 549 seconds, which contained an active engagement.

The model recognized changes in positioning, incoming
damage, a downed teammate, and a visible enemy knockdown.
It also identified that the player continued fighting
after receiving substantial damage.

However, the analysis failed to recognize the complete
engagement outcome. Although the player eliminated an
opposing squad, the generated report identified only
an individual knockdown.

The model consequently recommended retreating or
prioritizing a teammate revive without adequately
considering the successful elimination of the immediate
enemy threat.

This failure is particularly relevant to SPECTRE's
purpose: a tactical coach must understand the
development and resolution of an engagement before
judging the decisions made within it.

\section{Discussion}

The initial results reveal three issues that cannot
be resolved through HUD extraction alone.

\subsection{Recognition Is Not Game Understanding}

In one experiment, the VLM described a Mastiff
with five loaded rounds as critically low on ammunition,
despite five rounds representing a full magazine
in the tested configuration.

It also confused ammunition types and incorrectly
treated the Flatline's observed reserve ammunition
as a fixed maximum capacity.

These errors indicate that generic multimodal reasoning
is insufficient when tactical conclusions depend on
specialized weapon mechanics.

SPECTRE therefore requires an explicit, versioned
knowledge source describing weapon capacities,
ammunition types, legend abilities, and other
game-specific conventions.

\subsection{Missing Data Must Not Become False Evidence}

During combat analysis, the model repeatedly interpreted
\texttt{health: null} as evidence that the player was
downed. This contradicted the frames, which showed
continued weapon use and engagement.

The problem was not merely an inaccurate visual
prediction. It reflected a semantic mismatch between
the extraction schema and the model's interpretation
of that schema.

Future observations will distinguish unavailable
measurements, last-known values, and confirmed states.
The model must treat an absent measurement as
insufficient evidence rather than a negative event.

\subsection{Tactical Reasoning Requires Temporal Context}

A ten-second batch sampled at one frame per second
provides only ten visual observations. Rapid combat
events---including weapon switches, knockdowns,
eliminations, and movement between cover positions---may
occur entirely between these samples.

The missed squad wipe illustrates this limitation.
The model recognized an individual knockdown but lacked
sufficient temporal evidence, or failed to connect the
available evidence, to reconstruct the engagement outcome.

Another limitation arises at batch boundaries: a fight
may begin in one batch and end in the next. A short
natural-language summary of previous observations does
not necessarily preserve every relevant event.

The next architecture will therefore introduce
fight-aware sampling and explicit event memory.

\section{Future Work}

The next development stage prioritizes four extensions.

\textbf{Game-specific knowledge integration.}
A structured knowledge base will describe weapons,
legends, resource mechanics, and relevant HUD conventions.
Only entries relevant to the current match segment
will be supplied to the VLM.

\textbf{Fight-aware temporal sampling.}
The system will investigate higher sampling rates
during combat, initially targeting approximately
five frames per second. Surrounding context will
be retained to support interpretation of both
engagement initiation and outcome.

\textbf{Explicit state and event tracking.}
A separate state-management component will maintain
validated observations, last-known resource values,
and timestamped events such as damage, knockdowns,
and weapon changes.

\textbf{Domain adaptation and evaluation.}
Annotated Apex gameplay examples will be collected
to evaluate domain-specific errors and investigate
supervised fine-tuning after establishing a
knowledge-guided baseline.

Recognition of tactical abilities, inventory
contents, and audio events is also planned.
The longer-term objective remains event-driven,
deadline-aware coaching, in which late or
insufficiently supported advice is suppressed.

\section{Conclusion}

SPECTRE has reached its first functional milestone:
a complete offline pipeline that combines deterministic
HUD perception with sequential vision-language analysis.

The prototype demonstrates reliable integration of
video sampling, screen-state identification, selected
HUD detectors, temporal JSON observations, and
locally hosted multimodal inference.

Its experiments also establish an important architectural
lesson. Extracting visual facts is necessary but
insufficient for tactical coaching. Reliable feedback
requires game-specific knowledge, explicit treatment of
uncertainty, and temporal reasoning that connects
individual actions to engagement outcomes.

The next phase will address these limitations while
preserving the modular architecture established by
the current implementation.

\end{document}
